% ======================================================================
%  sec30.tex —— 第 30 节 整态射与有限态射
%  底本：原书 pp.161–163；原书标题：30. Integral and finite morphisms
% ======================================================================

\section{整态射与有限态射}\label{sec:30}

\subsection{定义}\label{sec:30.1}

回忆：若对任意 $b\in B$，都存在首一多项式 $p(x)\in A[x]$，使 $A$-同态 $A[x]\to B$，$x\mapsto b$ 把 $p(x)$ 送到零，则称环同态 $A\to B$ 是\emph{整的}。

回忆：若 $B$ 作为 $A$-模有限生成，则称环同态 $A\to B$ 是\emph{有限的}。已知有限同态都是整的。

\begin{definition}\label{def:30.1.1}
设 $f\colon X\to Y$ 是概形间的态射。若 $f$ 仿射，且对任意仿射开子概形 $U\subseteq Y$，环同态
\[
\cO_Y(U)\to\cO_X\bigl(f^{-1}(U)\bigr)
\]
都是整的（相应地，有限的），则称 $f$ 是\emph{整的}（相应地，\emph{有限的}）。
\end{definition}

\begin{warning}\label{warn:30.1.2}
整态射的概念并不是整概形的相对版本。
\end{warning}

\begin{example}\label{ex:30.1.3}
闭浸入是有限的。
\end{example}

\begin{lemma}\label{lem:30.1.4}
设 $\phi\colon A\to B$ 是环同态，$f\colon\Spec(B)\to\Spec(A)$ 是相应的概形态射。则 $f$ 整（相应地，有限）当且仅当 $\phi$ 整（相应地，有限）。
\end{lemma}

\begin{proof}
「仅当」显然。「当」由下述事实得到：若 $\phi$ 整（相应地，有限），则对任意环同态 $A\to A'$，基变换 $A'\to A'\otimes_AB$ 也整（相应地，有限）。
\end{proof}

\begin{lemma}\label{lem:30.1.5}
有限态射是整的。整且有限型的态射是有限的。
\end{lemma}


\begin{proof}
由交换代数中的相应结论得到。
\end{proof}

\begin{yizhu}
第二条即交换代数中的「整 + 有限型 $\Rightarrow$ 有限」：若 $B$ 由有限多个元素 $b_1,\ldots,b_n$ 作为 $A$-代数生成，则每个 $b_i$ 满足首一方程；由这些方程可取出有限生成 $A$-子模 $B'\subseteq B$，它含所有 $b_i$，故 $B'=B$。整性正是把「逐元素整」升级为「整体有限」的关键。
[来源：Stacks Project, Tag 02JJ（Lemma 10.36.5：$R\to S$ 整且有限型，当且仅当 $R\to S$ 有限），\url{https://stacks.math.columbia.edu/tag/02JJ}（2026-11-17 访问）。]
\end{yizhu}

下述结果的证明留给读者。

\begin{lemma}\label{lem:30.1.6}
考虑概形间的整（相应地，有限）态射的集合。有：
\begin{enumerate}[label=(\roman*),leftmargin=2.6em,itemsep=0.3ex]
  \item 设 $X\xrightarrow{f}Y\xrightarrow{g}Z$ 是态射链，且 $f$ 与 $g$ 都整（相应地，有限）。则 $g\circ f$ 整（相应地，有限）。
  \item 设 $X\xrightarrow{f}Y\xrightarrow{g}Z$ 是态射链，且 $g\circ f$ 整（相应地，有限）、$g$ 分离。则 $f$ 整（相应地，有限）。
  \item 整（相应地，有限）态射经基变换保持。
  \item 整性（相应地，有限性）对目标概形是局部的。
\end{enumerate}
\end{lemma}

\subsection{整态射与有限态射的纤维}\label{sec:30.2}

\begin{lemma}\label{lem:30.2.1}
设 $f\colon X\to Y$ 是整态射。则每点 $x\in X$ 都在纤维 $f^{-1}(f(x))$ 中闭。
\end{lemma}

\begin{proof}
记 $y=f(x)$。可沿 $\Spec(\kappa_y)\to Y$ 作基变换，从而设 $Y=\Spec(k)$，其中 $k$ 是域。此时断言由「$k$ 上整的整环是域」得到。
\end{proof}

\begin{yizhu}
所用事实是：若 $k$ 是域，$A$ 是整的 $k$-代数且为整环，则 $A$ 是域。证明用「$A$ 中非零元 $a$ 满足首一方程，故 $a$ 可逆」这一步，无需诺特假设。由此还可得：$k$ 上整的代数里每个素理想都是极大理想。
[来源：Stacks Project, Tag 00GS（Lemma 10.36.19：$k$-代数 $S$ 若整且为整环，则 $S$ 是域；且整 $k$-代数的每个素理想极大），\url{https://stacks.math.columbia.edu/tag/00GS}（2026-11-17 访问）。]
\end{yizhu}

\begin{lemma}\label{lem:30.2.2}
设 $f\colon X\to Y$ 是有限态射。则 $f$ 的每根纤维都是有限个点的有限不交并。
\end{lemma}

\begin{proof}
可化到 $Y=\Spec(k)$，其中 $k$ 是域。注意此时 $X$ 诺特，故只有有限多个不可约分支。由引理 \ref{lem:30.2.1}，任一不可约分支都是一点。这说明 $X$ 是有限个点的有限不交并。
\end{proof}

\begin{warning}\label{warn:30.2.3}
引理 \ref{lem:30.2.2} 的逆命题不成立。事实上，开浸入一般不是有限的。
\end{warning}

\begin{remark}\label{rem:30.2.4}
若态射 $f\colon X\to Y$ 有限型，且其每根纤维都是有限个点的有限不交并，则称 $f$ 是\emph{拟有限的}。这类态射将在《代数几何 II》中研究。
\end{remark}

\subsection{闭性}\label{sec:30.3}

\begin{proposition}\label{prop:30.3.1}
态射 $f\colon X\to Y$ 整，当且仅当它仿射且泛闭。
\end{proposition}

\begin{proof}
为证「仅当」，只需证整态射 $f\colon X\to Y$ 闭。由推论 \ref{cor:24.1.4}，只需证特化可沿 $f$ 提升。为此可设 $Y$ 从而 $X$ 仿射。此时断言由「整同态满足 going up」这一熟知事实得到（习题 \ref{exc:30.3.2}）。

为证「当」，可设 $Y$ 从而 $X$ 仿射。设 $\phi\colon A\to B$ 是对应于 $f$ 的环同态。设 $b\in B$ 是元素。只需证 $A[x]\to B$，$x\mapsto b$ 把某个首一多项式送到零。

考虑闭子概形 $Z:=\Spec(B[t]/(bt-1))\subseteq\A^1_X:=\Spec(B[t])$。记 $g$ 为复合
\[
g\colon Z\to\A^1_X\to\A^1_Y,
\]
并考虑其概形论像 $\overline{g(Z)}$（命题—定义 \ref{prop:18.5.4}）。按定义有
\[
\overline{g(Z)}\simeq\Spec(A[t]/I),
\]
其中
\[
I:=\ker\bigl(A[t]\to B[t]\to B[t]/(bt-1)\bigr).
\]
态射 $\A^1_X\to\A^1_Y$ 闭，因为 $f$ 泛闭。特别地，$g(Z)$ 是 $\A^1_Y$ 的闭子集。因 $g(Z)$ 在 $\overline{g(Z)}$ 中稠密，得 $Z\to\overline{g(Z)}$ 满。

考虑下述交换图
\[
\begin{tikzcd}[column sep=3.4em, row sep=2.6em]
\varnothing \arrow[r] \arrow[d] & \overline{g(Z)}\times_{\A^1_Y}Y \arrow[r] \arrow[d] & Y \arrow[d, "0"] \\
Z \arrow[r] & \overline{g(Z)} \arrow[r] & \A^1_Y,
\end{tikzcd}
\]
其中 $0\colon Y\to\A^1_Y$ 是对应于 $A[t]\to A$，$t\mapsto 0$ 的态射。注意右方块与外方块都是笛卡儿方块。这说明左方块也是笛卡儿的。因 $Z\to\overline{g(Z)}$ 满，得到
\[
\varnothing\simeq g(Z)\times_{\A^1_Y}Y\simeq\Spec\bigl(A[t]/(I+(t))\bigr).
\]
这说明 $1\in I+(t)$。故存在 $q(t)\in A[t]$，$q(0)=1$，使
\[
A[t]\to B[t]\to B[t]/(bt-1)
\]
把 $q(t)$ 送到零。这说明对 $d\ge\deg(q)$，
\[
A[t]\to B[t]\to B[t]/(t-b)
\]
把首一多项式 $t^dq(t^{-1})$ 送到零。
\end{proof}

\begin{yizhu}
「仅当」那一步的关键在于：整同态的 going up 恰好就是「特化可沿 $\Spec(B)\to\Spec(A)$ 提升」。而推论 \ref{cor:24.1.4} 判定闭性用的正是特化提升，两条语言在此对接。一般环同态的 going up 与特化提升等价，是可以逐素理想验证的拓扑翻译。
[来源：Stacks Project, Tag 00HW（Lemma 10.41.3：$R\to S$ 满足 going up，当且仅当特化可沿 $\Spec(S)\to\Spec(R)$ 提升），\url{https://stacks.math.columbia.edu/tag/00HW}（2026-11-17 访问）。]
\end{yizhu}

\begin{exercise}\label{exc:30.3.2}
设 $\phi\colon A\to B$ 是交换环间的整同态。设 $\mathfrak p\subseteq\mathfrak q\subseteq A$ 是 $A$ 的素理想，$\mathfrak p'\subseteq B$ 是 $B$ 的素理想且 $\phi^{-1}(\mathfrak p')=\mathfrak p$。则存在 $B$ 的素理想 $\mathfrak q'$，使 $\mathfrak p'\subseteq\mathfrak q'$ 且 $\phi^{-1}(\mathfrak q')=\mathfrak q$。
\end{exercise}

\begin{corollary}\label{cor:30.3.3}
态射 $f\colon X\to Y$ 有限，当且仅当它仿射且紧合。
\end{corollary}

\begin{corollary}\label{cor:30.3.4}
有限态射 $f\colon X\to Y$ 是射影的。
\end{corollary}


\begin{proof}
注意 $f_*\cO_X$ 是有限型拟凝聚 $\cO_Y$-模。故只需证
\[
X\to\PP_Y(f_*\cO_X)
\]
是闭浸入。问题对 $Y$ 是局部的，故可设 $Y$ 仿射。此时断言由定理 \ref{thm:29.1.5} 得到。
\end{proof}

\begin{yizhu}
两条推论合起来给出「有限 $\Leftrightarrow$ 仿射 + 紧合 $\Rightarrow$ 射影」的链条，这是整节的结构性结论。推论 \ref{cor:30.3.4} 的证明只用 $\cE=f_*\cO_X$ 有限型拟凝聚，再验证 $X\to\PP_Y(f_*\cO_X)$ 是闭浸入，即命题 \ref{prop:28.5.6} 的判据在 $\cL=\cO_X$ 时的特例。
[来源：Stacks Project, Tag 0B3I（Lemma 29.45.16：有限态射是射影的，经 $\sigma\colon X\to\mathbf P(f_*\cO_X)$ 为闭浸入），\url{https://stacks.math.columbia.edu/tag/0B3I}（2026-11-17 访问）。]
\end{yizhu}
